\documentclass[letterpaper, 10 pt, conference]{IEEEtran}
\IEEEoverridecommandlockouts
\overrideIEEEmargins

\usepackage{cite}
\usepackage{amsmath,amssymb,amsfonts}
\usepackage{algorithmic}
\usepackage{graphicx}
\usepackage{textcomp}
\usepackage{xcolor}
\usepackage{booktabs}
\usepackage{hyperref}

\title{\LARGE \bf
Autonomous Bimanual Soft-Body Slicing: Tactile-Conditioned Residual Reinforcement Learning under Contact Uncertainty
}

\author{Ahmed$^{*,1}$, Shahd$^{*,1}$, and Shan An$^{\dagger,1}$%
\thanks{$^{*}$Both authors contributed equally to this work.}%
\thanks{$^{1}$DEX-ROB Laboratory, School of Electrical and Automation Engineering, Tianjin University, Tianjin 300072, China.}%
\thanks{$^{\dagger}$Corresponding author: Prof. Shan An (email: {\tt shan.an@tju.edu.cn}).}%
}

\begin{document}

\maketitle
\thispagestyle{empty}
\pagestyle{empty}

\begin{abstract}
Robotic slicing of heterogeneous deformable food items, such as tomatoes, presents formidable challenges due to non-linear material stiffness, rupture mechanics, and severe contact uncertainty. Traditional pure kinematic trajectories frequently lead to excessive compressive deformation and pulp bursting before fracture occurs. In this paper, we propose a hierarchical vision-tactile residual reinforcement learning framework for autonomous bimanual slicing on dual 7-DoF manipulators (ARX AR5-L6) with a dexterous 5-finger fixturing hand (LinkerHand O6). Our approach decouples macro-sawing kinematics from high-frequency ($1\,\text{kHz}$) residual Cartesian impedance adaptation conditioned on 6-axis wrist wrench and fingertip tactile sensors. We validate our framework using GPU-accelerated PhysX 5 FEM simulation in NVIDIA Isaac Lab with extensive domain randomization over hyperelastic material properties. 
\end{abstract}

\section{Introduction}
\label{sec:intro}
Robotic manipulation of deformable objects (DOM) has gained substantial attention in surgical robotics, agricultural harvesting, and automated culinary preparation. Among these tasks, food slicing presents a unique scientific challenge: unlike rigid peg-in-hole assembly or non-destructive soft object grasping, cutting inherently induces irreversible structural topological fracture. 

Tomatoes embody a quintessential non-linear deformable material characterized by a high-toughness, flexible outer pericarp skin enclosing soft, fluid-filled locular gel and seeds. Applying classical rigid position control drives the knife downward regardless of surface deformation, creating immense internal hydrostatic pressure that results in crushing and fluid ejection long before skin puncture. Conversely, human chefs adaptively coordinate knife sawing frequency with subtle normal force adjustments based on continuous tactile and kinesthetic feedback.

In this work, we present an autonomous bimanual slicing framework that addresses these challenges through:
\begin{enumerate}
    \item A bimanual embodiment utilizing dual ARX AR5-L6 7-DoF arms with LinkerHand O6 adaptive fixturing and high-bandwidth instrumented cutting.
    \item A hierarchical residual reinforcement learning (RL) policy that superimposes micro-adjustments onto a nominal Cartesian sawing trajectory.
    \item Realistic GPU-accelerated PhysX 5 FEM deformable simulation in NVIDIA Isaac Lab with domain randomization over fracture energy and elastic moduli.
\end{enumerate}

\section{Related Work}
\label{sec:related_work}

\subsection{Robotic Slicing and Fracture Mechanics}
Early investigations into robotic cutting primarily focused on rigid food analogs or simplified fracture approximations. Sanchez et al.~\cite{sanchez2021tacblade} demonstrated the utility of tactile feedback during cutting of bio-materials using optical tactile sensors. However, most existing works rely on single-arm setups where the object is rigidly clamped in mechanical vices, precluding natural deformation dynamics encountered during human-like holding.

\subsection{Reinforcement Learning for Contact-Rich Manipulation}
Model-free reinforcement learning, specifically Proximal Policy Optimization (PPO)~\cite{schulman2017proximal}, has achieved remarkable results in dexterous in-hand manipulation and locomotion. Nevertheless, applying end-to-end RL to cutting leads to excessive sample complexity and unsafe exploration. Residual RL frameworks~\cite{johannink2019residual} mitigate this by learning corrective actions over nominal motion primitives, which we exploit to ensure contact stability.

\subsection{Simulation of Deformable Bodies and Sim-to-Real}
Recent advancements in GPU physics simulation, notably Isaac Gym and Isaac Lab~\cite{makoviychuk2021isaac}, enable real-time simulation of thousands of parallel environments. With PhysX 5, continuum mechanics and finite element models (FEM) can be executed natively on GPU tensor memory, facilitating systematic domain randomization for sim-to-real policy transfer.

\section{Problem Formulation}
\label{sec:problem_formulation}

\subsection{Fracture and Cutting Mechanics}
Let the deformable tomato be parameterized as a volumetric domain $\Omega \subset \mathbb{R}^3$ discretized by tetrahedral elements with material coordinates $\mathbf{X} \in \Omega_0$ and deformed coordinates $\mathbf{x}(\mathbf{X}, t) \in \Omega_t$. The deformation gradient is given by $\mathbf{F} = \frac{\partial \mathbf{x}}{\partial \mathbf{X}}$. The internal stress tensor $\boldsymbol{\sigma}$ follows a Neo-Hookean constitutive model parameterized by initial shear modulus $\mu$ and bulk modulus $\kappa$:
\begin{equation}
    \Psi(\mathbf{F}) = \frac{\mu}{2}(I_1 - 3 - 2\ln J) + \frac{\kappa}{2}(J - 1)^2
\end{equation}
where $J = \det(\mathbf{F})$ and $I_1 = \text{tr}(\mathbf{F}^T \mathbf{F})$. Rupture and cut propagation occur when the principal normal tensile stress $\sigma_{max}$ exceeds the critical rupture stress $\sigma_c$.

\subsection{Continuous Markov Decision Process}
The bimanual cutting task is modeled as a Markov Decision Process defined by the tuple $\langle \mathcal{S}, \mathcal{A}, \mathcal{P}, \mathcal{R}, \gamma \rangle$. 
The observation vector $\mathbf{o}_t \in \mathbb{R}^{33}$ captures the tool state, measured contact wrench $\mathbf{F}_{ee}, \mathbf{M}_{ee}$, LinkerHand tactile contact signals $\mathbf{s}_{tactile}$, cutting depth progress $z_{cut}$, and previous action $\mathbf{a}_{t-1}$.
The action space $\mathbf{a}_t \in \mathbb{R}^6$ specifies residual Cartesian velocities $\Delta \mathbf{v} \in \mathbb{R}^3$ and active Cartesian stiffness adaptations $\Delta \mathbf{K} \in \mathbb{R}^3$.

\section{Methodology}
\label{sec:method}

\subsection{Hierarchical Vision-Tactile Architecture}
Our architecture operates across two coupled temporal layers:
\begin{enumerate}
    \item \textbf{High-Level Kinematic Sawing Primitive ($30\,\text{Hz}$)}: Generates a baseline sinusoidal trajectory along the knife sawing axis $x$ and a nominal downward feed rate along $z$:
    \begin{equation}
        \mathbf{v}_{nom}(t) = [A \omega \cos(\omega t), 0, -v_{feed}]^T
    \end{equation}
    \item \textbf{Low-Level Residual Impedance Policy ($1\,\text{kHz}$)}: Evaluates policy $\pi_\theta(\mathbf{a}_t \mid \mathbf{o}_t)$ to output $\Delta \mathbf{v}_t$ and stiffness adjustment $\Delta \mathbf{K}_t$, producing the commanded end-effector velocity $\mathbf{v}_{cmd} = \mathbf{v}_{nom} + \Delta \mathbf{v}_t$.
\end{enumerate}

\subsection{Reward Function Design}
The composite scalar reward balances cutting throughput against crushing prevention:
\begin{equation}
    R_t = w_{p} R_{prog} + w_{e} \frac{|\dot{x}_{saw}|}{F_{z} + \epsilon} - w_{c} \mathbb{I}_{[\epsilon_{strain} > \epsilon_{crit}]} - w_{s} \|\mathbf{a}_t - \mathbf{a}_{t-1}\|^2
\end{equation}
where the force-efficiency ratio explicitly rewards shear sawing motion while penalizing excessive downward normal force $F_z$.

\section{Experimental Setup}
\label{sec:experiments}

\subsection{Simulation Environment in Isaac Lab}
We instantiate a bimanual cutting workstation in NVIDIA Isaac Lab atop Isaac Sim 6.1.0. The left arm houses an ARX AR5-L6 7-DoF manipulator with the LinkerHand O6 5-finger dexterous hand. The right arm controls the instrumented slicing knife. Deformable tomatoes are simulated using tetrahedral meshes solved via PhysX 5 FEM on GPU tensor cores. 

\subsection{Domain Randomization Parameters}
To bridge the sim-to-real gap, we apply domain randomization across key mechanical parameters:
\begin{itemize}
    \item Tomato Young's Modulus $E \sim \mathcal{U}(30\,\text{kPa}, 150\,\text{kPa})$.
    \item Surface friction coefficient $\mu \sim \mathcal{U}(0.2, 0.6)$.
    \item Puncture rupture stress $\sigma_c \sim \mathcal{U}(80\,\text{kPa}, 180\,\text{kPa})$.
    \item Random initial object pose offsets $\Delta \mathbf{p} \sim \mathcal{N}(\mathbf{0}, (5\,\text{mm})^2 \mathbf{I})$.
\end{itemize}

\subsection{Baseline Methods for Comparison}
We compare the proposed tactile residual RL policy against:
\begin{enumerate}
    \item \textbf{Kinematic Sawing}: Fixed open-loop trajectory without force feedback (\texttt{EXP-001}).
    \item \textbf{Fixed Impedance Control}: Classic Cartesian impedance with hand-tuned constant stiffness (\texttt{EXP-002}).
    \item \textbf{End-to-End Joint RL}: Monolithic 14-DoF PPO policy learning from scratch without kinematic primitive guidance (\texttt{EXP-004}).
\end{enumerate}

\section{Results and Ablation Studies}
\label{sec:results}

\subsection{Slicing Success Rate and Contact Regulation}
Quantitative benchmarks across $5$ random seeds for baseline and proposed methods will be evaluated upon completion of runs registered in \texttt{research\_state/experiment\_matrix.yaml}. 

\subsection{Deformation Suppression Analysis}
We analyze the maximum pre-fracture deformation ratio across firm, optimal, and overripe tomato categories to determine whether tactile impedance modulation prevents internal hydrostatic pressure spikes and pulp burst.

\subsection{Ablations and Sample Complexity}
We evaluate the learning curves of residual policy formulations against end-to-end joint and Cartesian policies, measuring the total environment interaction steps required to achieve stable cutting performance.

\subsection{Physical Sim-to-Real Transfer}
Zero-shot policy transfer onto the physical dual AR5-L6 testbed with LinkerHand O6 will benchmark real-world slice completion and boundary uniformity.

\section{Conclusion and Future Work}
\label{sec:conclusion}
We have presented an autonomous bimanual slicing framework for soft, deformable food items under severe contact uncertainty. By combining macro-sawing kinematic primitives with high-frequency tactile residual impedance adaptation in Isaac Lab, our system addresses the core challenge of crush prevention during fracture initiation. Future work will investigate vision-tactile multimodal foundation models and dynamic knife trajectory re-planning for irregular culinary geometries.


\bibliographystyle{IEEEtran}
\bibliography{references}

\end{document}
